\subsection{H4: individually removable is not jointly removable}

\HFourUnits{} test files met the size rule. more-itertools and dateutil contributed none. more-itertools keeps its tests in two files of 145 and 620 tests. dateutil's only file in range, \texttt{test\_imports.py}, has 27 tests, but only 8 of them record run-phase coverage, which leaves it below the 10-test minimum. tomlkit had one eligible file.

In every file, at least two tests were leave-one-out redundant, so all \HFourEligible{} files were evaluated. Removing the whole leave-one-out-redundant set together lost at least one kill in \HFourLosing{} of \HFourEligible{} files (\HFourShare{}, Wilson 95\% CI \HFourLo{} to \HFourHi{}). H4 is \textbf{\HFourSupported}. Across those files, group removal lost \HFourLostTotal{} of \HFourKilledTotal{} kills (Table~\ref{tab:h4}).

The extreme case is attrs' \texttt{test\_filters.py}: every one of its 33 tests is individually redundant, because every mutant it kills is killed by at least two tests, yet removing them all loses all 21 kills. Leave-one-out redundancy is a statement about one test given all the others, so it cannot be applied to several tests at once. This is the textbook result for set cover; the data shows it is not a corner case on real suites.

The greedy cover column shows how small a subset would preserve every observed kill: a median of \HFourCoverMedian{} of the file's tests (range \HFourCoverMin{} to \HFourCoverMax{}). That number is fitted to this mutant sample. It should be read as an upper bound on what a fast-feedback tier could drop, to be checked on independent mutants and historical bugs, not as a deletion list.

\begin{table}[t]
\centering\small
\begin{tabular}{lrrrrrrr}
\toprule
Project & Test file & Tests & Killed & LOO-redundant & Lost by group & Greedy cover & No kills \\
\midrule
attrs & test\_annotations.py & 47 & 31 & 44 & 11 & 6 & 0 \\
attrs & test\_filters.py & 33 & 21 & 33 & 21 & 4 & 0 \\
attrs & test\_funcs.py & 53 & 41 & 48 & 8 & 11 & 4 \\
toolz & toolz/test\_inspect\_args.py & 17 & 92 & 9 & 4 & 9 & 0 \\
toolz & toolz/test\_dicttoolz.py & 51 & 139 & 48 & 81 & 14 & 0 \\
toolz & toolz/test\_itertoolz.py & 51 & 138 & 18 & 3 & 34 & 1 \\
marshmallow & test\_utils.py & 20 & 45 & 13 & 1 & 8 & 5 \\
marshmallow & test\_options.py & 10 & 45 & 6 & 1 & 5 & 0 \\
marshmallow & test\_decorators.py & 44 & 84 & 41 & 41 & 8 & 0 \\
click & test\_info\_dict.py & 28 & 90 & 21 & 18 & 9 & 0 \\
click & test\_stream\_lifecycle.py & 23 & 82 & 19 & 2 & 5 & 4 \\
click & test\_testing.py & 43 & 89 & 35 & 8 & 11 & 3 \\
tomlkit & test\_parser.py & 34 & 80 & 28 & 12 & 8 & 0 \\
cachetools & test\_func.py & 36 & 112 & 24 & 6 & 15 & 0 \\
cachetools & test\_tlru.py & 34 & 126 & 25 & 7 & 13 & 0 \\
cachetools & test\_cache.py & 23 & 66 & 19 & 9 & 7 & 0 \\
\bottomrule
\end{tabular}
\caption{H4: kill matrices for sampled test files. \emph{LOO-redundant}: tests whose removal alone loses no kill. \emph{Lost by group}: kills lost when all of them are removed together. \emph{Greedy cover}: tests a greedy set cover keeps.}
\label{tab:h4}
\end{table}

